% BEGIN_REPORT::PROOF::WRP-0002::2026-10-07T16:14:00-07:00
\documentclass[11pt]{article}
\usepackage[margin=1in]{geometry}
\usepackage{amsmath,amssymb,amsthm,booktabs,hyperref}
\emergencystretch=3em
\newtheorem{theorem}{Theorem}
\newcommand{\SAT}{\mathrm{SAT}}
\title{WRP-0002: Certified Partial Portfolios for Exact Forward Wrappers}
\author{PROOF, ARC/WRAPPER research team}
\date{7 October 2026}
\begin{document}
\maketitle

\paragraph{Result and strict status.}
\textbf{CONDITIONAL promotion status; complete exact proof under the stated
hypotheses, with independent RED review pending.}
A fixed finite collection of partial forward wrappers can be absorbed into one
uniform inner procedure without a separate classifier, provided every
SUCCESS plan is exact, all failures halt within charged bounds, and the
declared success domains cover every input.  The combined fanout, child-size,
local-time, and live-space exponents are the coordinatewise maxima.  This
does not prove favorable coverage for general SAT.  That load-bearing arrow,
the ARC hitting premise, and P versus NP remain \textsc{open}.

\section*{Claim card}
\begin{description}
\item[] \textbf{Claim ID:} CERTIFIED-PARTIAL-PORTFOLIO1.
\item[] \textbf{Status before run:} the current WRAPPER spine stated the
fixed-portfolio idea informally; Drive records already proved finite-witness
portfolio absorption under all-safe/common-cap hypotheses.
\item[] \textbf{Exact quantifiers:} there exist a fixed $k$ and uniform
procedures $W_1,\ldots,W_k$ such that for every ordinary encoded CNF $F$, each
$W_j(F)$ halts with FAIL or an exact plan; every $F$ has at least one SUCCESS;
and all stated bounds hold for all sufficiently large input lengths.
\item[] \textbf{Model:} one deterministic random-access realization on the
fixed explicit SAT encoding.  SUCCESS/FAIL detection, discovery, discarded
attempts, output/generator representation, child access, combiner, and live
state are charged.
\item[] \textbf{SAME\_REALIZATION:} the sequential portfolio and its winning
plan are one machine.  Time from one member may not be paired with space from
another hypothetical execution.
\item[] \textbf{Weakest useful theorem:} prove the stated exact favorable
coverage using a fixed family of concrete partial mechanisms on all CNFs.
The present theorem then supplies one global inner interface automatically.
\end{description}

\section*{CORPUS REVIEWED}
The run executed the repository preflight and read root \texttt{main.tex},
\texttt{main2.tex}, \texttt{main3.tex}, current GitHub blobs of
\texttt{main4.tex} and \texttt{main5.tex}, \texttt{AGENTS.md},
\texttt{CODEX.md}, the full pnp-research skill and both protocol references,
the proof-spines protocol, current bibliography, WRP-0001 and its raw test,
the current ARC-0002 report, the wrapper-variant audit, proof-assembly audit,
double-wrapper implementation report, and repository searches for portfolios,
backdoors, coverage, failure flags, and wrapper variants.

The connected Drive sources were read rather than inferred from summaries:
SIM-1-0011 ``Finite-Witness Portfolio and Space Gate,'' SIM-2-0003
``Proof-Carrying Portfolio Gate,'' SIM-2-0007 ``Portfolio Absorption and
Dynamic-Witness Frontier,'' and ``Optimal Deterministic Backdoor-Cost
Recurrence.''  Their source URLs are
\url{https://docs.google.com/document/d/1Kbhlp0EK9zq0V7nNiEhfPcPKX0qS7hhNglpPbqAppdY},
\url{https://docs.google.com/document/d/1ko21wjx4WYmqtasf9aLbPtcRpkQHTt-fi3bDdOl4bEY},
\url{https://docs.google.com/document/d/1xuKhHO6Oh-k4ePd106HfeL2T5f4CmRaMtFIyDhGDqEI}, and
\url{https://docs.google.com/document/d/13fvjsNfRD-V9ZiembLbhg2kviF9Pg8NWq2Z7gXsDDOE}.

\section*{DUPLICATION / CONTRADICTION AUDIT}
The fixed-finite portfolio scheduling idea is not new.  SIM-1-0011 and
SIM-2-0007 already prove that a fixed favorable all-safe witness can be
absorbed by a charged scheduler, while warning that common small space and
favorable-member existence remain open.  SIM-2-0003 proves that merely one
correct first-finisher is unsafe.  This run adopts those results and does not
relabel them as a new general scheduling theorem.

The incremental result is a WRAPPER-specific specialization to partial
generators with atomic FAIL/SUCCESS, exact combination plans, and explicit
fanout/size/acquisition exponents.  It converts the earlier informal
``portfolio still needs coverage'' sentence into the exact sufficient
interface and shows that no learned or semantic selector is needed once the
finite cover is proved.  No contradiction with the current root manuscripts
was found.

\section*{The exact interface theorem}
\begin{theorem}[CERTIFIED-PARTIAL-PORTFOLIO1]
Fix $k<\infty$.  On a SAT instance $F$ of encoded length $M$, each uniform
$W_j$ halts with FAIL or a plan
$(C_j;G_{j,1},\ldots,G_{j,q_j})$.  Suppose:
\begin{enumerate}
\item each SUCCESS plan satisfies
\[
\SAT(F)=C_j(\SAT(G_{j,1}),\ldots,\SAT(G_{j,q_j}));
\]
\item for every $F$, at least one $W_j$ succeeds;
\item success and failure, including discarded search, take
$M^{c_j+o(1)}$ time and $M^{s_j+o(1)}$ live space, without a hidden
current-scale SAT call;
\item a successful plan has $q_j\le M^{a_j+o(1)}$ children, each of encoded
length at most $M^{b_j+o(1)}$ for $b_j<1$, with representation, generation,
access and combination charged.
\end{enumerate}
Then sequential trial until the first SUCCESS is one exact uniform inner
procedure with
\[
a=\max_j a_j,
\quad b=\max_j b_j<1,
\quad c=\max_j c_j,
\quad s=\max_j s_j.
\]
\end{theorem}
\begin{proof}
Coverage produces a SUCCESS, and every SUCCESS is semantically exact, so the
first returned plan is exact.  Since $k$ is constant, the sum of all attempt
times has exponent $\max_j c_j$.  Sequential execution discards failed state,
so peak space has exponent $\max_j s_j$.  Only the winning plan is retained,
giving the fanout and child-size maxima.  The maximum of finitely many
$o(1)$ functions remains $o(1)$.  Explicit output writing or implicit
generator state cannot be omitted because the hypotheses charge the actual
representation.
\end{proof}

Recursive closure follows only if every produced child is again in the same
covered ordinary-SAT domain.  The theorem does not excuse a retained failed
cache, an unsafe member, a one-sided SAT-only validator, an input-dependent
number of candidates, or a dynamic certificate whose acquisition is not in
$c_j,s_j$.

\section*{Primary literature and transfer audit}
M{\"a}hlmann, Siebertz and Vigny define recursive component backdoors and
derive tractability for bounded recursive depth \cite{maehlmann2021}.
Dreier, Ordyniak and Szeider provide FPT approximation and linear-time SAT
algorithms for bounded backdoor depth into Horn and Krom
\cite{dreier2022}.  These are \textbf{THEOREM TRANSFERS} only for their
declared structural domains: such backdoors yield exact partial wrappers.
They do not establish bounded depth, favorable encoded shrink, or cheap
acquisition on every CNF.

SATzilla selects among complementary solvers per instance
\cite{xu2008}.  Classification: \textbf{ANALOGY/CANDIDATE}.  Its empirical
hardness models and distributional objectives do not provide theorem-level
worst-case coverage or a joint-resource certificate.  The missing bridge is
exact favorable coverage in the present random-access/ordinary-CNF model.
The post-quantum collision, Schur-square and binding toolkit does not match
this scheduling/interface obstacle; no algebraic object or cryptographic
hardness assumption is imported.

\section*{Bounded executable model}
The new runner
\path{experiments/proof_spines/certified_partial_portfolio_test.py} uses
seed 20261007, a 25-second cooperative alarm, and standard exact mechanisms:
constant detection, Horn forward chaining, 2-SAT implication SCCs,
variable-disjoint component AND, and exhaustive discovery of a strong Horn
backdoor of size at most two followed by an OR plan.  Production wrappers make
zero SAT-oracle calls.  Brute force is used only as the independent bounded
test oracle.

The preserved run completed in 0.053136 seconds.  All 320 generated formulas
(80 from each declared family) produced correct plans.  There were 1,072
wrapper attempts, including 752 charged failures and 568 charged candidate
restrictions.  Ninety-six random order permutations preserved the answer.
Winning counts were 123 Horn, 40 2-CNF, 79 components, and 78 size-two Horn
backdoors.  Raw output is
\path{results/raw/proof_spines/2026-10-07_161400_certified_partial_portfolio.json};
the source digest is
\texttt{e6f51e58befc7b0a7ea570423d9c5bb208b049d201abd87450c8992af0f75d36}.

\section*{Hostile controls and failed coverage}
\begin{description}
\item[] \textbf{Outside coverage:} a connected seven-variable cycle of
positive width-three clauses caused all five members to return FAIL.  This is
not a hardness result; it proves the executable portfolio is not a general-SAT
inner wrapper.
\item[] \textbf{Unsafe first finisher:} inserting a fast false-answer member
before an exact member gives a wrong result on a satisfiable clause.  Thus
``some member is correct'' is insufficient.
\item[] \textbf{Independent product:} four disconnected width-three clauses
were split exactly into four children.  The easy win says nothing about
connected general inputs.
\item[] \textbf{Padding:} sixty-four fresh tautological clauses were removed
by normalization and supplied no contraction.
\item[] \textbf{Growing portfolio:} a synthetic $k(M)=M=4096$ last-winner
case requires 4096 sequential attempts; constant-factor absorption no longer
applies without a separate exponent and memory ledger.
\item[] \textbf{Source carry/current scale:} production makes zero calls to
the hypothetical SAT decider.  Adding such a call would retain its unknown
exponent.  Children carrying the full source would also need charged access
and would not automatically satisfy the ordinary-SAT shrink bound.
\item[] \textbf{Oracle/algebraic replacement:} the theorem survives giving
all members a common oracle and survives algebraic extension access.  It is
therefore infrastructure, not the nonrelativizing or nonalgebrizing hinge.
\end{description}

\section*{WHAT CHANGED RELATIVE TO ROOT MAIN}
Root \texttt{main.tex} was not edited: it is the unrelated dashboard with
separate ownership regions.  Root \texttt{main5.tex} now contains the exact
partial-portfolio theorem, proof, transfer limits and WRP-0002 status entry.
Root \texttt{main4.tex} receives a concise ARC review recording no change to
the cofinal-seed OPEN premise.  The bibliography adds verified primary entries
for recursive/depth backdoors and SATzilla.

\section*{ADOPTED / REJECTED / RETIRED PRIOR COMPONENTS}
\begin{description}
\item[] \textbf{Adopted:} finite-witness portfolio absorption; common-space
gate; exactness-by-construction; recursive backdoors as exact restricted
mechanisms; charged backdoor recurrence; sequential discard of failed state.
\item[] \textbf{Rejected:} empirical winner prediction as proof of coverage;
one correct member among unsafe first finishers; one-sided YES validation;
automatic small-space preservation; treating a growing candidate family as a
constant portfolio.
\item[] \textbf{Retired:} none beyond the already retired unsafe-portfolio
formulation.  The new bounded portfolio itself is explicitly not promoted to
general SAT.
\end{description}

\section*{STALE OR OVERSTATED CLAIMS FOUND}
No live manuscript claim required downgrade.  The older informal statement
that a portfolio needs favorable coverage and aggregate memory was correct
but incomplete; it is now sharpened.  Historical reports that treat selection
of one fixed safe member as a separate asymptotic obstacle are superseded by
SIM-1-0011/SIM-2-0007 under their common-cap hypotheses.

\section*{Open arrow and peer handoff}
The remaining wrapper target is now exact:
\[
\boxed{\text{fixed finite exact partial mechanisms}
\quad+\quad\text{favorable charged coverage of every CNF}.}
\]
PROOF should next test whether component, recursive-backdoor, frontier-width,
and source-backed mechanisms admit one common fixed-power $(a,b,c,s)$ envelope
on a larger explicitly defined domain.  RED should minimize a CNF outside any
proposed union and test retained failed state, UNSAT exactness, padding, source
carry, and dynamic-member growth.  RESEARCH should seek a primary theorem that
proves worst-case coverage rather than instance-distribution success.  Without
that coverage theorem, the joint WRAPPER interface remains \textsc{open}.

\bibliographystyle{plain}
\bibliography{spines}
% END_OF_VERIFIED_REPORT::PROOF::WRP-0002
\end{document}
